\section{Optimal release cardinality}

Output cardinality is a property of the stationary release design that attains the optimized treatment-effect information. Although the optimization built from the fourteen patterns in \cref{obj:synth_8,obj:synth_4} uses all candidate rays, an efficient design can always be chosen with at most five active outputs. The following result establishes this bound throughout the interior trial model.

\begin{theoremv}{thm:five-output-upper-bound}[Five-output optimal staircase channel]
Let \(\theta=(\mu_0,\mu_1)^\top\), \(p\), and \(\varepsilon\) satisfy:
\begin{itemize}
\item \textbf{(Interior means.)} \(0<\mu_0<1\) and \(0<\mu_1<1\).
\item \textbf{(Assignment probability.)} \(0<p<1\).
\item \textbf{(Privacy parameter.)} \(\varepsilon>0\).
\end{itemize}
With the feasible staircase weights and profiled objective of \cref{obj:def:staircase-saddle}, there exists \(\alpha\in\mathcal A_\varepsilon\) such that
\[
J^*(\theta,p,\varepsilon)
=
\inf_{t\in\mathbb R}F_\theta(\alpha,t),
\qquad
\bigl|\{S\in\mathcal S:\alpha_S\ne0\}\bigr|\le5,
\qquad
\kappa(Q_\alpha)\le5.
\]
Here the active-output cardinality of the associated fourteen-symbol staircase channel is
\[
\kappa(Q_\alpha)
=
\bigl|\{S\in\mathcal S:(P_\theta Q_\alpha)(\{S\})\ne0\}\bigr|.
\]
\end{theoremv}

The fourteen patterns provide a common representation for the optimization, while the active weights determine the alphabet used by a particular release. \Cref{obj:thm:five-output-upper-bound} therefore gives a uniform bound on the number of output symbols needed to attain the optimized profiled information. Its practical content is an existence guarantee: at each admissible parameter point, the efficient stationary experiment admits a release with at most five active symbols.

To identify where particular supports attain the optimum, we use certificates that combine channel normalization, minimization over the profiling coordinate, and dual optimality for the weights. The fixed pattern enumeration in \cref{obj:synth_8} gives the support indices a common meaning across parameter points.

\begin{definitionv}{synth_5}[Staircase support certificates]
Fix \(\theta\in(0,1)^2\), \(p\in(0,1)\), and \(\varepsilon>0\), and use the staircase rays, masses, and derivative vectors defined above. For a proposed active set \(A\subseteq\mathcal S\), a support certificate consists of weights \(\alpha\in\mathbb R_+^{14}\), a profiling coordinate \(t\in\mathbb R\), and a dual vector \(\eta\in\mathbb R^4\), with \(\alpha_S>0\) for \(S\in A\) and \(\alpha_S=0\) for \(S\notin A\). Write
\[
f_S(t)=\frac{(g_S^\top v(t))^2}{h_S(\theta)}.
\]
We say that primal normalization holds when
\[
B\alpha=\mathbf1,
\qquad\text{equivalently}\qquad
\sum_{S\in\mathcal S}\alpha_S b_S=\mathbf1,
\]
where \(B=(b_{S_1},\ldots,b_{S_{14}})\).

We say that profiling stationarity holds at \((\alpha,t)\) when
\[
\frac{\partial}{\partial t}F_\theta(\alpha,t)
=\sum_{S\in\mathcal S}\alpha_S f_S'(t)
=2\sum_{S\in\mathcal S}\alpha_S
\frac{(g_S^\top v(t))(g_S^\top(1,1)^\top)}{h_S(\theta)}
=0.
\]
For feasible \(\alpha\), the function \(F_\theta(\alpha,\cdot)\) is a differentiable convex quadratic, so this condition is equivalent to \(t\in\arg\min_{u\in\mathbb R}F_\theta(\alpha,u)\).

The dual vector \(\eta\) corresponds to the normalization constraint in the linear program
\[
\max_{\alpha\ge0:\,B\alpha=\mathbf1}
\sum_{S\in\mathcal S}\alpha_S f_S(t);
\]
its dual feasibility conditions are \(b_S^\top\eta\ge f_S(t)\) for every \(S\in\mathcal S\). We say that active dual equalities hold when
\[
b_S^\top\eta=f_S(t)\qquad(S\in A),
\]
and that strict inactive dual inequalities hold when
\[
b_S^\top\eta>f_S(t)\qquad(S\in\mathcal S\setminus A).
\]
The active sets for the three-output and five-output certificates are respectively
\[
A=\{S_1,S_6,S_8\},
\qquad
A=\{S_3,S_4,S_6,S_8,S_9\},
\]
using the fixed pattern enumeration. The five-output certificate additionally requires pairwise distinct values of
\[
\rho_S(\theta,t)=\frac{g_S^\top v(t)}{h_S(\theta)}
\qquad(S\in A).
\]
\end{definitionv}

Normalization makes the weights a feasible channel, and profiling stationarity selects a minimizing coordinate for its directional information. The dual vector \leanref{sym:\eta}{\ensuremath{\eta}} bounds each pattern's directional contribution \leanref{sym:f_S(t)}{\ensuremath{f_S(t)}}; equality on active patterns identifies where the bound is attained. Strict inequalities on inactive patterns distinguish the proposed support within this optimization. The additional separation of projected scores in the five-output certificate supports the cardinality conclusion for arbitrary attaining stationary channels. These conditions define the following parameter regions.

\begin{definitionv}{def:support-regions}[Support regions $\mathcal R_3$ and $\mathcal R_5$]
The region \(\mathcal R_3\) is the parameter set for which, using the fixed pattern enumeration \(S_1,\ldots,S_{14}\), the following conditions hold with active set \(\{S_1,S_6,S_8\}\):
\begin{itemize}
\item \textbf{(Normalization.)} Primal normalization holds.
\item \textbf{(Stationarity.)} Profiling stationarity holds.
\item \textbf{(Active patterns.)} The active dual equalities hold.
\item \textbf{(Inactive patterns.)} The inactive dual inequalities are strict.
\end{itemize}
The region \(\mathcal R_5\) is the analogous parameter set with active set \(\{S_3,S_4,S_6,S_8,S_9\}\), additionally satisfying
\[
\rho_S(\theta,t)\ne\rho_T(\theta,t)
\qquad
\text{for all distinct }S,T\in\{S_3,S_4,S_6,S_8,S_9\}.
\]
\end{definitionv}

The regions \leanref{sym:\mathcal R_3}{\ensuremath{\mathcal R_3}} and \leanref{sym:\mathcal R_5}{\ensuremath{\mathcal R_5}}, defined in \cref{obj:def:support-regions}, describe parameters admitting the respective support certificates. Their coordinates are ordered as treatment-assignment probability, control-arm mean, treatment-arm mean, and privacy parameter. The five-output region makes the global existence bound sharp on an open parameter set and gives explicit examples at equal arm means.

\begin{theoremv}{thm:five-output-region}[The five-output region]
The five-pattern certificate and score-separation region \(\mathcal R_5\), associated with the support \(\{S_3,S_4,S_6,S_8,S_9\}\), is open in \(\mathbb R^4\) and nonempty. Write \(J^*(\theta,p,\varepsilon)\) and \(V^*(\theta,p,\varepsilon)\) for the optimized profiled information and its reciprocal from \cref{obj:def:staircase-saddle}.

For every \(0<p<1/2\),
\[
(p,1/2,1/2,\log 3)\in\mathcal R_5.
\]
At \(\theta=(1/2,1/2)^\top\) and \(\varepsilon=\log 3\), there is a unique feasible staircase-weight vector \(\alpha\) attaining the optimized profiled information:
\[
J^*(\theta,p,\log 3)
=
\inf_{t\in\mathbb R}F_\theta(\alpha,t),
\qquad
\{S\in\mathcal S:\alpha_S\ne0\}
=
\{S_3,S_4,S_6,S_8,S_9\}.
\]
Here feasibility means that the weights are nonnegative and their weighted staircase rays sum to one at each of the four input records. Uniqueness means that every feasible weight vector \(\beta\) satisfying
\[
J^*(\theta,p,\log 3)
=
\inf_{t\in\mathbb R}F_\theta(\beta,t)
\]
equals \(\alpha\).

For every \(1/2<p<1\),
\[
(1-p,1/2,1/2,\log 3)\in\mathcal R_5.
\]
At \(\theta=(1/2,1/2)^\top\) and the original assignment probability \(p\), there is likewise a unique feasible staircase-weight vector \(\alpha\) satisfying
\[
J^*(\theta,p,\log 3)
=
\inf_{t\in\mathbb R}F_\theta(\alpha,t).
\]
Its support is the image of \(\{S_3,S_4,S_6,S_8,S_9\}\) under exchange of the treatment labels.

In each of these two cases, every stationary \((\log 3)\)-LDP channel \(Q\) from the four-record alphabet to an arbitrary measurable output space that attains contrast variance \(V^*(\theta,p,\log 3)\) satisfies \(\kappa(Q)\ge5\). More generally, for every
\[
(p,\mu_0,\mu_1,\varepsilon)\in\mathcal R_5,
\qquad
\theta=(\mu_0,\mu_1)^\top,
\]
every stationary \(\varepsilon\)-LDP channel \(Q\) attaining contrast variance \(V^*(\theta,p,\varepsilon)\) satisfies
\[
\kappa(Q)\ge5.
\]
A stationary \(\varepsilon\)-LDP channel is a Markov kernel satisfying
\[
Q(x,A)\le e^\varepsilon Q(x',A)
\]
for every measurable output event \(A\) and every pair of input records \(x,x'\). Attainment means that the extended nonnegative contrast variance of the output Fisher-information matrix \(I_\theta(Q)\) equals the nonnegative part of \(V^*(\theta,p,\varepsilon)\). Explicitly, this contrast variance is
\[
\begin{cases}
\displaystyle
\max\!\left\{
\inf_{v:\,I_\theta(Q)v=c}c^\top v,\;0
\right\},
& c\in\operatorname{range}(I_\theta(Q)),\\[4pt]
+\infty,
& c\notin\operatorname{range}(I_\theta(Q)),
\end{cases}
\]
where \(c\) is the treatment-effect contrast vector. The output-cardinality convention is
\[
\kappa(Q)=
\begin{cases}
\#\{z\in\mathcal Z:(P_\theta Q)(\{z\})\ne0\},
& \mathcal Z\ \text{finite},\\
+\infty,
& \mathcal Z\ \text{infinite},
\end{cases}
\]
where \(P_\theta Q\) is the channel's output law.
\end{theoremv}

Five outputs here reflect a requirement of efficient stationary release design. Throughout the certified five-output region, the lower bound applies to every attaining stationary channel under the stated output-cardinality convention. Together with \cref{obj:thm:five-output-upper-bound}, it makes five the sharp cardinality bound on this region. At equal arm means and privacy parameter \(\log 3\), the theorem also identifies a unique optimizing staircase-weight vector for each assignment probability below one half; exchanging treatment labels gives the corresponding characterization above one half. The uniqueness conclusion concerns weights in the fixed staircase representation, while the cardinality conclusion concerns the full class of attaining stationary channels.

A smaller optimal support occurs at the same privacy parameter when both arm means are low. The next result provides an exact three-pattern optimizer and its information value.

\begin{theoremv}{thm:three-output-region}[Three-output region and optimizer]
The three-pattern certificate region \(\mathcal R_3\) is open in \(\mathbb R^4\) and contains
\[
(p,\mu_0,\mu_1,\varepsilon)
=\left(\frac{3}{10},\frac{1}{10},\frac{1}{10},\log 3\right).
\]
At this point, with \(\theta=(1/10,1/10)^\top\), there exists a feasible staircase-weight vector \(\alpha\), in the optimization of \cref{obj:def:staircase-saddle}, whose support and weights satisfy
\[
\{S\in\mathcal S:\alpha_S\ne0\}=\{S_1,S_6,S_8\},
\qquad
\alpha_{S_1}=\alpha_{S_6}=\alpha_{S_8}=\frac15.
\]
This vector attains the optimized profiled information:
\[
J^*\left(\theta,\frac{3}{10},\log 3\right)
=\inf_{t\in\mathbb R}F_\theta(\alpha,t),
\]
where the profiled objective is evaluated at \(p=3/10\) and \(\varepsilon=\log 3\). Moreover, every feasible staircase-weight vector \(\beta\) at \(\varepsilon=\log 3\) satisfying
\[
J^*\left(\theta,\frac{3}{10},\log 3\right)
=\inf_{t\in\mathbb R}F_\theta(\beta,t)
\]
equals \(\alpha\). The optimized information and its reciprocal variance, as defined in \cref{obj:def:staircase-saddle}, are
\[
J^*\left(\theta,\frac{3}{10},\log 3\right)=\frac{441}{3994},
\qquad
V^*\left(\theta,\frac{3}{10},\log 3\right)=\frac{3994}{441}.
\]
\end{theoremv}

At this parameter point, three active patterns suffice to achieve the optimized information, and their equal weights give an explicit benchmark for the finite optimization. The rational information and variance values permit a direct comparison with numerical solutions. Openness of the certificate region places this example within a neighborhood of certified parameters. The following result connects the local staircase supports to exact minimum output cardinalities among arbitrary attaining stationary channels.

\begin{theoremv}{thm:support-transition}[Local support and cardinality transition]
Fix the treatment-assignment probability \(p=3/10\) and privacy parameter \(\varepsilon=\log 3\). There exist disjoint open sets \(U_5,U_3\subseteq\mathbb R^2\) containing \((1/2,1/2)\) and \((1/10,1/10)\), respectively, with the following properties:
\begin{itemize}
\item For every \(\theta=(\mu_0,\mu_1)^\top\in U_5\), there exists a unique feasible staircase-weight vector \(\alpha\in\mathcal A_\varepsilon\) satisfying
\[
J^*(\theta,p,\varepsilon)
=
\inf_{t\in\mathbb R}F_\theta(\alpha,t).
\]
Its active support is \(\{S_3,S_4,S_6,S_8,S_9\}\). Moreover, a feasible staircase channel attains contrast variance \(\max\{V^*(\theta,p,\varepsilon),0\}\) with exactly five active outputs, and every stationary \(\varepsilon\)-LDP channel on any measurable output space attaining this contrast variance has at least five active outputs.
\item For every \(\theta=(\mu_0,\mu_1)^\top\in U_3\), there exists a unique feasible staircase-weight vector \(\alpha\in\mathcal A_\varepsilon\) satisfying
\[
J^*(\theta,p,\varepsilon)
=
\inf_{t\in\mathbb R}F_\theta(\alpha,t).
\]
Its active support is \(\{S_1,S_6,S_8\}\). Moreover, a feasible staircase channel attains contrast variance \(\max\{V^*(\theta,p,\varepsilon),0\}\) with exactly three active outputs, and every stationary \(\varepsilon\)-LDP channel on any measurable output space attaining this contrast variance has at least three active outputs.
\end{itemize}
In each case, uniqueness is among all feasible staircase-weight vectors satisfying the displayed equality.
\end{theoremv}

Holding assignment and privacy fixed, the minimum cardinality changes from five near arm means \((1/2,1/2)\) to three near \((1/10,1/10)\). Each conclusion persists on an open neighborhood, so the comparison captures a local dependence of efficient release design on outcome prevalence. In both neighborhoods, uniqueness selects the optimizer within the staircase family, and the matching universal lower bound establishes the minimum cardinality for all attaining stationary channels. The certificate algebra and the projected-score arguments establishing these cardinality bounds are presented in \cref{sec:proofs}.

Balanced assignment and complementary arm means give a further interpretable special case of the general optimization. Under this symmetry, binary randomized response applied to the treatment--outcome sign attains the optimized contrast variance.

\begin{theoremv}{thm:balanced-reduction}[Balanced randomized response reduction]
Let \(\theta=(\mu_0,\mu_1)^\top\) be the arm-mean parameter and \(\tau=\mu_1-\mu_0\) the treatment-effect contrast. Suppose:
\begin{itemize}
\item \textbf{(Interior means.)} \(\theta\) satisfies \cref{obj:ass:interior-means}.
\item \textbf{(Balance.)} \(p=1/2\) and \(\mu_0+\mu_1=1\).
\item \textbf{(Privacy parameter.)} \(\varepsilon>0\).
\end{itemize}
Let \(Q\) be binary randomized response applied to \((2W_i-1)(2Y_i-1)\): it releases the input sign with probability
\[
\frac{e^\varepsilon}{e^\varepsilon+1}
\]
and the opposite sign with probability
\[
\frac{1}{e^\varepsilon+1}.
\]
Write \(V^*(\theta,p,\varepsilon)=J^*(\theta,p,\varepsilon)^{-1}\) for the reciprocal optimized nuisance-profiled information. Then
\[
V^*(\theta,1/2,\varepsilon)
=
\left(\frac{\cosh(\varepsilon/2)}{\sinh(\varepsilon/2)}\right)^2-\tau^2.
\]
The channel \(Q\) is a stationary Markov channel satisfying, for every measurable output event and every pair of input records,
\[
Q(x_j,A)\le e^\varepsilon Q(x_{j'},A).
\]
Moreover, let \(I_\theta(Q)\) be the Fisher-information matrix of its output experiment at \(p=1/2\), and let \(c=(-1,1)^\top\). With contrast variance defined as \(c^\top I_\theta(Q)^+c\) when \(c\in\operatorname{range}(I_\theta(Q))\), and as \(+\infty\) otherwise, this channel attains
\[
c^\top I_\theta(Q)^+c=V^*(\theta,1/2,\varepsilon),
\qquad
c\in\operatorname{range}(I_\theta(Q)).
\]
\end{theoremv}

The input sign records whether the treatment and outcome signs agree, aligning the binary release with the treatment-effect contrast. Under balanced assignment and complementary arm means, this release retains the information needed for efficient contrast estimation: the range condition makes estimability explicit, and its contrast variance equals the optimized bound. Both release probabilities are positive, giving exactly two active outputs. The closed-form variance therefore supplies a simple benchmark for the general optimization under the stated symmetry, alongside the three- and five-output designs characterized above.
